\section{Marco teórico}

\subsection{Bases teóricas}
El OBR se apoya en la detección de objetos de una etapa, que localiza y clasifica
la celda Braille completa en una sola pasada. La cuantización INT8 reduce el
tamaño y la latencia del modelo, a costa de una posible pérdida de precisión que
debe acotarse con un criterio de no inferioridad. La evaluación del desempeño
combina métricas de precisión, concordancia y eficiencia, y se complementa con
medidas de carga cognitiva y usabilidad.

\subsection{Definición de términos}
\begin{description}
  \item[Braille Grado 1.] Escritura táctil en la que cada carácter se representa
  por una celda de seis puntos, sin contracciones.
  \item[OBR.] Reconocimiento óptico de Braille: conversión de una imagen de texto
  en relieve a texto digital.
  \item[YOLO.] Familia de detectores de objetos de una etapa que predice cajas y
  clases en una sola pasada.
  \item[Cuantización INT8.] Reducción de los pesos y activaciones de un modelo a
  enteros de 8 bits para disminuir memoria y latencia.
\end{description}
